% Part: first-order-logic
% Chapter: syntax-and-semantics
% Section: formation-sequences

\documentclass[../../../include/open-logic-section]{subfiles}

\begin{document}

\olfileid{pl}{syn}{fseq}

\olsection{Vormingsrijen}

Het definiëren van !!{formula}s met een inductieve definitie en de
aanvullende techniek om eigenschappen van !!{formula}s door inductie te
bewijzen, vormen een elegante en efficiënte benadering. Het kan echter
ook nuttig zijn om de constructie van !!{formula}s van onderaf en stap
voor stap te beschouwen. Dat doen we hier met het begrip
\emph{vormingsrij}.

\begin{defn}[Vormingsrijen voor formules]
\ollabel{defn:fseq-frm}
Een eindige rij $\tuple{!A_0,\dotsc,!A_n}$ van tekenreeksen uit de
taal~$\Lang L_0$ is een \emph{vormingsrij} voor $!A$ als
$!A \ident !A_n$ en voor alle $i \leq n$ geldt: $!A_i$ is een
atomaire formule, of er bestaan $j,k < i$ waarvoor een van de volgende
gevallen geldt:
\begin{enumerate}
    \tagitem{prvNot}{$!A_i \ident \lnot !A_j$.}{}%
    \tagitem{prvAnd}{$!A_i \ident (!A_j \land !A_k)$.}{}%
    \tagitem{prvOr}{$!A_i \ident (!A_j \lor !A_k)$.}{}%
    \tagitem{prvIf}{$!A_i \ident (!A_j \lif !A_k)$.}{}%
    \tagitem{prvIff}{$!A_i \ident (!A_j \liff !A_k)$.}{}%
\end{enumerate}
\end{defn}

\begin{ex}
\[
\tuple{
    \Obj p_0,
    \Obj p_1,
    (\Obj p_1 \land \Obj p_0),
    \lnot (\Obj p_1 \land \Obj p_0)
}
\]
is een vormingsrij van
$\lnot (\Obj p_1 \land \Obj p_0)$, evenals
\[
\tuple{
    \Obj p_0,
    \Obj p_1,
    \Obj p_0,
    (\Obj p_1 \land \Obj p_0),
    (\Obj p_0 \lif \Obj p_1),
    \lnot (\Obj p_1 \land \Obj p_0)
}.
\]
%
Zoals uit het tweede voorbeeld blijkt, mogen vormingsrijen ``rommel''
bevatten: formules die overbodig zijn of niet aan de constructie
bijdragen.
\end{ex}

\begin{prop}
\ollabel{prop:formed}
Elke !!{formula}~$!A$ in~$\Frm[L_0]$ heeft een vormingsrij.
\end{prop}

\begin{proof}
Stel dat $!A$ atomair is. Dan is de rij~$\tuple{!A}$ een vormingsrij
voor~$!A$.
%
Stel nu dat $!B$ en~$!C$ respectievelijk de vormingsrijen
$\tuple{!B_0,\dotsc,!B_n}$ en $\tuple{!C_0,\dotsc,!C_m}$ hebben.
%
\begin{enumerate}
  \tagitem{prvNot}{Als $!A \ident \lnot !B$, dan is
      $\tuple{!B_0,\dotsc,!B_n,\lnot !B_n}$
      een vormingsrij voor~$!A$.}{}
  \tagitem{prvAnd}{Als $!A \ident (!B \land !C)$, dan is
      $\tuple{!B_0,\dotsc,!B_n,!C_0,\dotsc,!C_m,(!B_n \land !C_m)}$
      een vormingsrij voor~$!A$.}{}
  \tagitem{prvOr}{Als $!A \ident (!B \lor !C)$, dan is
      $\tuple{!B_0,\dotsc,!B_n,!C_0,\dotsc,!C_m,(!B_n \lor !C_m)}$
      een vormingsrij voor~$!A$.}{}
  \tagitem{prvIf}{Als $!A \ident (!B \lif !C)$, dan is
      $\tuple{!B_0,\dotsc,!B_n,!C_0,\dotsc,!C_m,(!B_n \lif !C_m)}$
      een vormingsrij voor~$!A$.}{}
  \tagitem{prvIff}{Als $!A \ident (!B \liff !C)$, dan is
      $\tuple{!B_0,\dotsc,!B_n,!C_0,\dotsc,!C_m,(!B_n \liff !C_m)}$
      een vormingsrij voor~$!A$.}{}
\end{enumerate}
Volgens het inductieprincipe voor !!{formula}s heeft elke !!{formula}
een vormingsrij.
\end{proof}

We kunnen ook het omgekeerde bewijzen. Dit is belangrijk, omdat het
aantoont dat onze twee definities van formules equivalent zijn: zij
leveren dezelfde resultaten op. Het betekent ook dat we stellingen
over formules kunnen bewijzen door gewone inductie naar de lengte van
vormingsrijen.

\begin{lem}
\ollabel{lem:fseq-init}
Stel dat $\tuple{!A_0,\dotsc,!A_n}$ een vormingsrij voor~$!A_n$ is en
dat $k \leq n$. Dan is $\tuple{!A_0,\dotsc,!A_k}$ een vormingsrij
voor~$!A_k$.
\end{lem}

\begin{proof}
Opgave.
\end{proof}

\begin{thm}
\ollabel{thm:fseq-frm-equiv}
$\Frm[L_0]$ is de verzameling van alle tekenreeksen in de
taal~$\Lang L_0$ die een vormingsrij hebben.
\end{thm}

\begin{proof}
Zij $F$ de verzameling van alle tekenreeksen in de taal~$\Lang L_0$
die een vormingsrij hebben. In \olref[pl][syn][fseq]{prop:formed}
hebben we gezien dat $\Frm[L_0] \subseteq F$. We bewijzen nu de
omgekeerde inclusie.

Stel dat $!A$ een vormingsrij $\tuple{!A_0,\dotsc,!A_n}$ heeft. We
bewijzen dat $!A \in \Frm[L_0]$ met sterke inductie naar~$n$. Onze
inductiehypothese luidt dat elke tekenreeks met een vormingsrij waarvan
de laatste index $m < n$ is, tot~$\Frm[L_0]$ behoort. Volgens de
definitie van een vormingsrij is $!A_n$ atomair, of er bestaan
$j,k < n$ waarvoor een van de volgende gevallen geldt:
\begin{enumerate}
    \tagitem{prvNot}{$!A_n \ident \lnot !A_j$.}{}%
    \tagitem{prvAnd}{$!A_n \ident (!A_j \land !A_k)$.}{}%
    \tagitem{prvOr}{$!A_n \ident (!A_j \lor !A_k)$.}{}%
    \tagitem{prvIf}{$!A_n \ident (!A_j \lif !A_k)$.}{}%
    \tagitem{prvIff}{$!A_n \ident (!A_j \liff !A_k)$.}{}%
\end{enumerate}
We onderscheiden nu gevallen. Als $!A_n$ atomair is, dan is
$!A_n \in \Frm[L_0]$. Stel daarentegen dat $!A \ident
(!A_j \land !A_k)$. Volgens \olref[pl][syn][fseq]{lem:fseq-init} zijn
$\tuple{!A_0,\dotsc,!A_j}$ en $\tuple{!A_0,\dotsc,!A_k}$ vormingsrijen
voor respectievelijk $!A_j$ en~$!A_k$. Omdat het echte beginstukken
van de vormingsrij voor~$!A$ zijn, zijn hun laatste indices kleiner
dan~$n$. Volgens de inductiehypothese behoren $!A_j$ en~$!A_k$ dus
  tot~$\Frm[L_0]$. Uit de definitie van !!a{formula} volgt dan dat dit
  dus ook voor $(!A_j \land !A_k)$ geldt. De overige gevallen
volgen uit analoge redeneringen.
\end{proof}

\end{document}
